%%%%%%%%%%  scan idx 619  |  folio 606  %%%%%%%%%%

En
effet on peut évidemment supposer \(S\) et \(T\) affines. Prenons un recouvrement affine \(\{V_i\}\) de \(X_T\). Soient \(g\in\Gamma(X_T;{O}_{X'_T}^{*})\)~; \(g_i\) un relèvement de \(g|V_i\) dans \(\Gamma(V_i;{O}_{X_T}^{*})\). Alors \(\delta_T^*(g)\) est égal à la classe du \(1\)-cocycle \(\delta_{ij}^*=(g_ig_j^{-1}-1)\)~; \(\delta_T(g)\), à celle de \(\delta_{ij}=(g_i-g_j)=g_j\delta_{ij}^*=g\,\delta_{ij}^*\), donc on a
\[
  \delta_T(g)=g\,\delta_T^*(g)
\]
en tenant compte de l'opération de \(\Gamma(X_T,\underline{O}_{X'_T})\) sur \(I_T\), donc sur \(H^1(X_T,I_T)\).

D'après ii), on a \(\Gamma(X_T;{O}_{X'_T})=\Gamma(X;{O}_{X'})\otimes_S\Gamma(T,{O}_T)\) et donc \(g=\sum g_\beta\otimes t_\beta\) pour
\(g_\beta\in\Gamma(X,{O}_{X'})\) et
\(t_\beta\in\Gamma(T,{O}_T)\). Or par (v) \((f_*{O}_X)_{\red}\longrightarrow(f_*{O}_{X'})_{\red}\) est surjectif~; par suite \(g_\beta=v(h_\beta)+n_\beta\) pour
\(h_\beta\in\Gamma(X;{O}_X)\), \(v:\Gamma(X,{O}_X)\to
\Gamma(X,{O}_{X'})\) et
\(n_\beta\in\operatorname{nilrad}(\Gamma(X,{O}_{X'}))\).

D'après i),
\((n_\beta\otimes t_\beta)I_T=0\). Par suite \(\delta_T(g)=h\,\delta_T^*(g)\) où \(h=\sum h_\beta\otimes t_\beta\in\Gamma(X_T,\underline{O}_{X_T})\) d'où \(\delta_T^*(g)=h^{-1}\delta_T(g)=\delta_T(h^{-1}g)\), donc \(\im\delta_T^*\subset\im\delta_T\).

Finalement soit \(g=\sum g_\beta\otimes t_\beta\in
\Gamma(X_T,{O}_{X'_T})\). Soit \(g_\beta=v(h_\beta)+n_\beta\) où
\(n_\beta\in\operatorname{nilrad}(\Gamma(X,{O}_{X'}))\); on a alors \(\delta_T(g)=\delta_T\bigl(\sum n_\beta\otimes t_\beta\bigr)=\delta_T\bigl(1+\sum n_\beta\otimes t_\beta\bigr)=\delta_T^*\bigl(1+\sum n_\beta\otimes t_\beta\bigr)\), d'où \(\im\delta_T\subset\im\delta_T^*\).

En vertu de 3.2, iii) et iv), on a une suite exacte fonctorielle en \(T\)~:
\[
  0\longrightarrow \mathbb V(E^\vee)(T)\longrightarrow
  R^1(f_T)_*({O}_{X_T}^{*})
  \longrightarrow
  R^1(f_T)_*({O}_{X'_T}^{*}) \longrightarrow \mathbb V(Q)(T),
\]
qui par passage aux faisceaux f.p.p.f. associés donne une suite exacte
\[
  0\longrightarrow \mathbb V(E^\vee)\longrightarrow\underline{\Pic}_{X/S}
  \longrightarrow\underline{\Pic}_{X'/S}\xrightarrow{u}\mathbb V(Q).
\]

Si \(P=\Ker(u)\), le monomorphisme \(P\hookrightarrow\underline{\Pic}_{X'/S}\) s'obtient par changement de base à partir de la section unité \(S\hookrightarrow\mathbb V(Q)\)~; comme \(\mathbb V(Q)\) est représentable et séparé, il résulte que \(P\hookrightarrow\Pic_{X'/S}\) est représentable par une
